% !TEX root = ../main.tex

\chapter{پیاده‌سازی سامانه}

% =========================================================
% 5-1
% =========================================================
\section{مقدمه}

در این فصل، جزئیات پیاده‌سازی سامانه گفت‌وگوی صوتی فارسی
تشریح می‌شود.

سامانه بر اساس معماری ماژولار معرفی‌شده در فصل قبل توسعه
یافته است و هسته اصلی آن شامل سه بخش
\lr{ASR}،
\lr{LLM}
و
\lr{TTS}
است.

هدف اصلی در پیاده‌سازی، فراهم‌کردن بستری بوده است که در آن
مدل‌های مختلف بتوانند با کمترین تغییر در سایر بخش‌های برنامه
جایگزین شوند.

در کنار این سه مؤلفه، بخش‌هایی برای دریافت صوت، پردازش متن،
مدیریت مکالمه، ثبت زمان اجرا و مدیریت خطا نیز در نظر گرفته
شده‌اند.


% =========================================================
% 5-2
% =========================================================
\section{محیط توسعه}

پیاده‌سازی اصلی پروژه با زبان
\lr{Python}
انجام شده است.

دلیل انتخاب \lr{Python}، پشتیبانی گسترده کتابخانه‌های یادگیری
ماشین، پردازش صوت و مدل‌های زبانی از این زبان بوده است.

بسته به مدل مورد استفاده، کتابخانه‌هایی مانند
\lr{PyTorch}،
\lr{Transformers}،
\lr{NeMo}،
\lr{Vosk}
و ابزارهای مرتبط با پردازش صوت در سامانه استفاده می‌شوند.

مدل‌ها تا حد امکان به‌صورت محلی بارگذاری و اجرا می‌شوند تا
وابستگی خط لوله اصلی به سرویس‌های خارجی کاهش یابد.


% =========================================================
% 5-3
% =========================================================
\section{ساختار نرم‌افزاری پروژه}

ساختار نرم‌افزار به‌صورت ماژولار طراحی شده است.
به‌طور مفهومی می‌توان بخش‌های پروژه را به شکل زیر در نظر گرفت:

\begin{latin}\begin{verbatim}
project/
|
|-- asr/
|   |-- base.py
|   |-- models/
|
|-- llm/
|   |-- base.py
|   |-- models/
|
|-- tts/
|   |-- base.py
|   |-- models/
|
|-- audio/
|-- conversation/
|-- evaluation/
|-- utils/
|-- main.py
\end{verbatim}\end{latin}

هدف از این ساختار، جداکردن مسئولیت‌های مختلف سامانه از یکدیگر
است.

برای مثال، کد مربوط به یک مدل \lr{ASR} نباید مستقیماً به جزئیات
مدل زبانی وابسته باشد.
ارتباط میان مؤلفه‌ها تنها از طریق ورودی و خروجی تعریف‌شده
انجام می‌شود.


% =========================================================
% 5-4
% =========================================================
\section{پیاده‌سازی ماژول \lr{ASR}}

ماژول \lr{ASR} مسئول دریافت فایل یا آرایه صوتی و تولید متن فارسی
است.

رابط کلی این ماژول به‌صورت مفهومی شامل عملیاتی مشابه زیر است:

\begin{latin}\begin{verbatim}
load_model()
transcribe(audio)
unload_model()
\end{verbatim}\end{latin}

تابع
\lr{load\_model}
مدل و پردازنده‌های مورد نیاز را در حافظه بارگذاری می‌کند.

تابع
\lr{transcribe}
صوت ورودی را دریافت کرده و خروجی متنی استاندارد تولید می‌کند.

در پروژه، چند خانواده مختلف \lr{ASR} از طریق همین منطق مشترک
ارزیابی شده‌اند، از جمله:

\begin{itemize}

    \item \lr{wav2vec2-large-xlsr-persian-v3}
    \item \lr{Whisper Large-v3}
    \item \lr{vosk-model-fa-0.42}
    \item \lr{stt\_fa\_fastconformer\_hybrid\_large}

\end{itemize}

تفاوت در نحوه بارگذاری و پیش‌پردازش هر مدل داخل همان ماژول
مدیریت می‌شود تا خروجی همه مدل‌ها برای بخش بعدی یک متن فارسی
باشد.


% =========================================================
% 5-5
% =========================================================
\section{پیاده‌سازی مدل زبانی}

پس از تولید متن توسط \lr{ASR}، متن در اختیار ماژول مدل زبانی قرار
می‌گیرد.

رابط اصلی این بخش را می‌توان به شکل زیر نمایش داد:

\begin{latin}\begin{verbatim}
load_model()
generate(prompt, history)
unload_model()
\end{verbatim}\end{latin}

ورودی مدل شامل پیام فعلی کاربر و در صورت فعال‌بودن مدیریت
مکالمه، بخشی از تاریخچه پیام‌های قبلی است.

در ارزیابی‌های پروژه چند مدل مختلف استفاده شده‌اند:

\begin{itemize}

    \item \lr{gemma-2-2b-it}
    \item \lr{gemma-2-9b-it}
    \item \lr{Qwen2.5-3B-Instruct}
    \item \lr{Qwen2.5-7B-Instruct}
    \item \lr{Dorna2-Llama3.1-8B-Instruct}

\end{itemize}

برای مقایسه منصفانه، همه این مدل‌ها در ارزیابی‌های اصلی با
کوانتیزه‌سازی چهاربیتی \lr{NF4} و تنظیمات یکسان اجرا شدند. دلیل
محدودکردن گزینه‌ها به مدل‌های 2 تا 9 میلیارد پارامتری نیز این بود
که نسخه کوانتیزه‌شده آن‌ها بتواند به‌طور کامل در حافظه یک \lr{GPU}
محلی بارگذاری شود. این شرط، امکان اجرای مستقل از سرویس ابری و
مقایسه مدل‌ها روی سخت‌افزار یکسان را فراهم کرد.

کوانتیزه‌سازی در اینجا عمدتاً دقت ذخیره‌سازی وزن‌ها را کاهش می‌دهد؛
در نتیجه نباید مقدار نظری حافظه وزن‌ها را با کل حافظه مصرفی فرآیند
یکسان دانست. فعال‌سازی‌ها، حافظه نهان \lr{KV}، طول بافت و سربار
کتابخانه استنتاج همچنان بر بیشینه مصرف حافظه اثر می‌گذارند. به همین
دلیل، قابلیت استقرار بر اساس مصرف واقعی در زمان اجرا بررسی شد، نه
فقط برآورد حاصل از تعداد پارامترها.

برای هر مدل، قالب گفت‌وگو و توکن‌ساز متناظر با همان مدل
استفاده می‌شود.

خروجی نهایی این بخش به شکل یک رشته متنی استاندارد در اختیار
ماژول \lr{TTS} قرار می‌گیرد.


% =========================================================
% 5-6
% =========================================================
\section{مدیریت تاریخچه مکالمه}

برای پشتیبانی از مکالمه چندمرحله‌ای، پیام‌های کاربر و پاسخ‌های
سامانه در یک ساختار تاریخچه نگهداری می‌شوند.

به‌صورت ساده:

\begin{latin}\begin{verbatim}
history = [
    {"role": "user", "content": "..."},
    {"role": "assistant", "content": "..."}
]
\end{verbatim}\end{latin}

در هر نوبت، پیام جدید به تاریخچه افزوده شده و متن مناسب برای
ورودی مدل زبانی تولید می‌شود.

با افزایش طول مکالمه، تعداد توکن‌ها نیز افزایش می‌یابد.
بنابراین در طراحی نهایی باید محدودیتی برای تعداد پیام‌های
نگهداری‌شده یا طول کل تاریخچه اعمال شود.


% =========================================================
% 5-7
% =========================================================
\section{پیاده‌سازی ماژول \lr{TTS}}

ماژول \lr{TTS} پاسخ متنی مدل زبانی را دریافت کرده و فایل یا آرایه
صوتی تولید می‌کند.

رابط کلی آن به شکل زیر در نظر گرفته شده است:

\begin{latin}\begin{verbatim}
load_model()
synthesize(text)
unload_model()
\end{verbatim}\end{latin}

مدل‌های بررسی‌شده در این بخش عبارت‌اند از:

\begin{itemize}

    \item \lr{Kamtera\_VITS}
    \item \lr{VITS\_MMS\_FAS}
    \item \lr{PIPER\_TTS\_FAS}

\end{itemize}

پیش از ارسال متن به مدل \lr{TTS}، در صورت نیاز پردازش‌هایی مانند
حذف کاراکترهای نامناسب، مدیریت علائم و نرمال‌سازی اولیه انجام
می‌شود.

خروجی مدل در قالب صوت ذخیره شده یا مستقیماً برای پخش ارسال
می‌شود.


% =========================================================
% 5-8
% =========================================================
\section{اتصال ماژول‌ها}

هسته اصلی برنامه، خروجی هر ماژول را به ورودی مرحله بعدی
منتقل می‌کند.

منطق ساده‌شده سامانه به شکل زیر است:

\begin{latin}\begin{verbatim}
audio = record_audio()

text = asr.transcribe(audio)

response = llm.generate(
    text,
    history
)

speech = tts.synthesize(response)

play_audio(speech)
\end{verbatim}\end{latin}

این ساختار ساده باعث می‌شود جریان اصلی برنامه مستقل از مدل
خاص مورد استفاده باقی بماند.

در نتیجه، تغییر مدل \lr{ASR} یا \lr{TTS} نیازی به بازنویسی منطق اصلی
مکالمه ندارد.


% =========================================================
% 5-9
% =========================================================
\section{رابط خط فرمان}

برای آزمایش و استفاده اولیه سامانه، یک رابط خط فرمان
\LTRfootnote{Command-Line Interface (CLI)}
در نظر گرفته شده است.

در این رابط، کاربر می‌تواند یک نوبت مکالمه را آغاز کند،
متن تشخیص‌داده‌شده را مشاهده کرده و پاسخ مدل زبانی را قبل
از پخش صوت بررسی کند.

نمایش خروجی‌های میانی در مرحله توسعه اهمیت زیادی دارد، زیرا
کمک می‌کند مشخص شود خطای نهایی از کدام مؤلفه ناشی شده است.

به‌طور نمونه، رابط می‌تواند اطلاعات زیر را نمایش دهد:

\begin{latin}\begin{verbatim}
\lr{ASR}: ...
\lr{LLM}: ...
\lr{TTS}: generated
\lr{Latency}: ...
\end{verbatim}\end{latin}

این رابط همچنین برای انجام تست‌های انتهابه‌انتها مورد استفاده
قرار می‌گیرد.


% =========================================================
% 5-10
% =========================================================
\section{ثبت زمان اجرا}

برای اندازه‌گیری زمان اجرای هر ماژول، قبل و بعد از اجرای
عملیات اصلی زمان ثبت می‌شود.

به‌صورت مفهومی:

\begin{latin}\begin{verbatim}
start = timer()
result = module.run(...)
end = timer()

latency = end - start
\end{verbatim}\end{latin}

این اندازه‌گیری برای \lr{ASR}، \lr{LLM} و \lr{TTS} به‌صورت جداگانه انجام
می‌شود.

سپس تأخیر کلی یک نوبت مکالمه نیز از مجموع مراحل یا اندازه‌گیری
مستقیم کل خط لوله محاسبه می‌شود.

این داده‌ها در فصل آزمایش‌ها برای مقایسه مدل‌ها و مشخص‌کردن
گلوگاه سامانه استفاده خواهند شد.


% =========================================================
% 5-11
% =========================================================
\section{مدیریت \lr{CPU} و \lr{GPU}}

مدل‌های مختلف از نظر نیاز سخت‌افزاری تفاوت زیادی دارند.

برخی مدل‌ها برای دستیابی به سرعت مناسب روی \lr{GPU} اجرا می‌شوند،
در حالی که برخی مدل‌های سبک‌تر مانند بعضی مدل‌های \lr{Vosk}
قابلیت اجرای مناسب روی \lr{CPU} را نیز دارند.

بنابراین در زمان بارگذاری مدل، دستگاه محاسباتی با توجه به
مدل و منابع در دسترس انتخاب می‌شود.

یکی از نکات مهم در اجرای چند مدل، آزادسازی حافظه پس از پایان
استفاده از یک مؤلفه در صورت کمبود منابع است.

این موضوع به‌ویژه زمانی اهمیت دارد که \lr{ASR}، \lr{LLM} و \lr{TTS} همگی
برای استفاده از یک \lr{GPU} رقابت کنند.


% =========================================================
% 5-12
% =========================================================
\section{مدیریت خطا}

برای جلوگیری از توقف کامل برنامه، اجرای ماژول‌ها باید در برابر
خطاهای رایج کنترل شود.

نمونه خطاهای ممکن شامل موارد زیر هستند:

\begin{itemize}

    \item عدم دریافت صوت معتبر؛
    \item شکست در بارگذاری مدل؛
    \item کمبود حافظه \lr{GPU}؛
    \item تولید متن خالی توسط \lr{ASR}؛
    \item شکست تولید پاسخ مدل زبانی؛
    \item و خطا در تولید یا پخش صوت.
\end{itemize}

در صورت بروز خطا، اطلاعات مربوط به آن ثبت شده و در صورت امکان
سامانه بدون بسته‌شدن کامل به وضعیت قابل استفاده بازمی‌گردد.


% =========================================================
% 5-13
% =========================================================
\section{وضعیت مؤلفه‌های تکمیلی}

مسیر اصلی

\begin{center}
\lr{ASR}
$\longrightarrow$
\lr{LLM}
$\longrightarrow$
\lr{TTS}
\end{center}

در نسخه فعلی پروژه پیاده‌سازی شده است.

در کنار این خط لوله، چند مؤلفه برای تکمیل سامانه در نظر گرفته
شده‌اند که شامل:

\begin{itemize}

    \item تشخیص فعالیت صوتی؛
    \item مقاومت و پردازش نویز؛
    \item بهبود علائم نگارشی؛
    \item مدیریت کامل‌تر تاریخچه مکالمه؛
    \item بهبود مدیریت خطا؛
    \item و اندازه‌گیری دقیق تأخیر انتهابه‌انتها
\end{itemize}

هستند.

این موارد به‌صورت جداگانه توسعه و سپس از نظر تأثیر بر عملکرد
کل سامانه ارزیابی خواهند شد.


% =========================================================
% 5-14
% =========================================================
\section{جمع‌بندی}

در این فصل، نحوه پیاده‌سازی معماری ماژولار پروژه تشریح شد.

سه ماژول اصلی \lr{ASR}، \lr{LLM} و \lr{TTS} با رابط‌های مشخص در نظر گرفته
شده‌اند و جریان اصلی برنامه خروجی هر بخش را به مرحله بعدی
منتقل می‌کند.

همچنین مدیریت تاریخچه، رابط خط فرمان، ثبت زمان اجرا، مدیریت
منابع و کنترل خطا به‌عنوان بخش‌های مکمل پیاده‌سازی معرفی شدند.

این ساختار امکان مقایسه مدل‌های مختلف در شرایط یکسان و سپس
یکپارچه‌سازی مدل‌های منتخب در خط لوله نهایی را فراهم می‌کند.

در فصل بعد، روش انجام آزمایش‌ها، مجموعه‌داده‌ها، معیارهای
ارزیابی و نتایج کمی مدل‌های مختلف ارائه خواهند شد.
